\documentclass[12pt,a4paper]{article}

% ── Packages ──────────────────────────────────────────────────────────────────
\usepackage[margin=2.5cm]{geometry}
\usepackage{graphicx}
\usepackage{booktabs}
\usepackage{listings}
\usepackage{xcolor}
\usepackage{hyperref}
\usepackage{parskip}
\usepackage{titlesec}
\usepackage{enumitem}
\usepackage{float}
\usepackage{amsmath}
\usepackage{microtype}

% ── Hyperref config ───────────────────────────────────────────────────────────
\hypersetup{
    colorlinks=true,
    linkcolor=black,
    urlcolor=blue,
    citecolor=black,
}

% ── Section styling ───────────────────────────────────────────────────────────
\titleformat{\section}{\large\bfseries}{}{0em}{}[\titlerule]
\titleformat{\subsection}{\normalsize\bfseries}{}{0em}{}

% ── SQL code listing style ────────────────────────────────────────────────────
\lstdefinestyle{sql}{
    language=SQL,
    basicstyle=\ttfamily\small,
    keywordstyle=\color{blue}\bfseries,
    stringstyle=\color{red!70!black},
    commentstyle=\color{gray}\itshape,
    backgroundcolor=\color{gray!8},
    frame=single,
    rulecolor=\color{gray!40},
    framesep=4pt,
    breaklines=true,
    showstringspaces=false,
    tabsize=2,
    captionpos=b,
}

\lstdefinestyle{prompt}{
    basicstyle=\ttfamily\footnotesize,
    backgroundcolor=\color{blue!4},
    frame=single,
    rulecolor=\color{blue!30},
    framesep=4pt,
    breaklines=true,
    showstringspaces=false,
}

% ── Document ──────────────────────────────────────────────────────────────────
\begin{document}

% ── Title page ────────────────────────────────────────────────────────────────
\begin{titlepage}
    \centering
    \vspace*{2cm}
    {\LARGE\bfseries Natural Language to SQL:\\[0.4em]
    Building a Text-to-SQL Agent for the Chicago Crime Dataset\par}
    \vspace{1.5cm}
    {\large Applied Data Science --- NLP Project\par}
    \vspace{0.5cm}
    {\large\itshape Course Report\par}
    \vspace{2cm}
    {\normalsize March 2026\par}
    \vfill
    \hrule
    \vspace{0.5cm}
    {\small
    \textbf{Dataset:} Chicago Crime Dataset (2020 -- Present)\\
    \textbf{Database:} PostgreSQL (\texttt{chicago\_crime})\\
    \textbf{Models evaluated:} \texttt{llama3.1:8b} (general-purpose) \ \textbullet\
    \texttt{mannix/defog-llama3-sqlcoder-8b} (fine-tuned)
    }
    \vspace{0.5cm}
    \hrule
\end{titlepage}

\tableofcontents
\newpage

% ─────────────────────────────────────────────────────────────────────────────
\section{Project Overview}
% ─────────────────────────────────────────────────────────────────────────────

The goal of this project was to build a system that accepts natural-language questions
about crime in Chicago and returns accurate, data-backed answers by automatically
generating and executing SQL queries against a live PostgreSQL database.

The Chicago Crime dataset contains records from 2020 to the present, covering attributes
such as crime type, date, district, arrest status, and geographic coordinates.
Rather than hard-coding queries, we wanted an interface that any analyst could use
without SQL knowledge.

The core research question became: \textit{how well can an off-the-shelf or fine-tuned
language model generate correct SQL on its own, with no retrieval step and no training
on our specific schema?}

% ─────────────────────────────────────────────────────────────────────────────
\section{What Went Well}
% ─────────────────────────────────────────────────────────────────────────────

\subsection{RAG was not needed}

Early in the project we planned to use \textbf{Retrieval-Augmented Generation (RAG)}
as a core part of the pipeline.  The intuition was that the database schema and any
relevant data samples would first be retrieved from a vector store based on the
user's question, and only then passed to the LLM.  This added a significant layer of
complexity: we would have needed to embed the schema, maintain a vector index,
implement a retriever, and manage chunking strategies for the DDL.

We discovered that the schema for our database is compact enough to fit entirely
inside the LLM's context window --- the \texttt{chicago\_crime} table has a
manageable number of columns and a single-table schema. This meant we could
\textbf{inject the full DDL directly into the prompt} on every call without any
retrieval step.

This simplification had several benefits:
\begin{itemize}[leftmargin=*, itemsep=2pt]
    \item \textbf{Reduced infrastructure:} no vector database, no embedding model,
          no retrieval latency.
    \item \textbf{Deterministic context:} the model always sees all columns, avoiding
          the risk of a retriever omitting a relevant field.
    \item \textbf{Faster iteration:} the pipeline became a straightforward
          \emph{prompt $\rightarrow$ LLM $\rightarrow$ SQL $\rightarrow$ execute}
          chain.
\end{itemize}

\subsection{LangGraph agent worked end-to-end}

The multi-step LangGraph agent (using \texttt{llama3.1:8b}) successfully completed
the full pipeline: listing tables $\to$ fetching schema $\to$ generating SQL $\to$
reviewing SQL $\to$ executing $\to$ summarising.  This validated the architecture
and provided a strong baseline to compare against.

\subsection{Robust fallback handling}

Because Ollama models sometimes emit plain JSON or raw SQL text instead of structured
tool calls, a three-pass fallback strategy was implemented in \texttt{sql\_agent.py}.
This made the LangGraph agent resilient to model output variation without requiring
a fine-tuned model.

% ─────────────────────────────────────────────────────────────────────────────
\section{Model Comparison: \texttt{llama3.1:8b} vs.\ SQLCoder}
% ─────────────────────────────────────────────────────────────────────────────

We evaluated two models under the same database and the same natural-language questions.
Table~\ref{tab:model_comparison} summarises the key differences.

\begin{table}[H]
\centering
\caption{High-level comparison of the two evaluated models}
\label{tab:model_comparison}
\renewcommand{\arraystretch}{1.4}
\begin{tabular}{lll}
\toprule
\textbf{Property} & \textbf{llama3.1:8b} & \textbf{SQLCoder (defog-llama3-8b)} \\
\midrule
Training objective     & General-purpose chat       & Fine-tuned for text-to-SQL \\
Tool / function calls  & Supported (unreliable)     & Not supported \\
How SQL is requested   & Tool call (\texttt{sql\_db\_query}) & Completion from prompt \\
Schema delivery        & Via schema tool in loop    & Full DDL injected in prompt \\
Pipeline complexity    & Multi-step LangGraph graph & Single-pass prompt chain \\
Identifier quoting     & Requires explicit prompting & Learned from fine-tuning \\
Hallucinated columns   & Occasional                 & Rare \\
Execution success rate & Moderate (needs fallbacks) & Higher out-of-the-box \\
\bottomrule
\end{tabular}
\end{table}

\subsection{Why a general-purpose model underperforms on SQL}

\texttt{llama3.1:8b} is a strong general reasoning model, but it was not trained
specifically for SQL generation.  Several practical problems surfaced during evaluation:

\begin{enumerate}[leftmargin=*, itemsep=3pt]
    \item \textbf{Unreliable tool invocation.}  The model occasionally ignored the
          \texttt{tool\_choice="any"} instruction and printed a JSON object in plain
          text rather than making a structured tool call.  Custom fallback parsing was
          required.

    \item \textbf{Identifier quoting.}  Column names that contain spaces or mixed case
          (e.g.\ \texttt{Primary Type}, \texttt{Location Description}) were sometimes
          left unquoted or wrapped in backticks instead of PostgreSQL-style double
          quotes, causing \texttt{syntax error} exceptions.

    \item \textbf{Table name hallucination.}  Without explicit instruction, the model
          occasionally invented table names (\texttt{crimes}, \texttt{crime\_data})
          that do not exist in the database.

    \item \textbf{Query complexity.}  For multi-step aggregation questions the model
          sometimes produced overly complex subqueries or omitted a GROUP BY clause.
\end{enumerate}

SQLCoder, by contrast, was fine-tuned on a corpus of (schema, question, SQL) triples.
It has internalised the PostgreSQL quoting conventions, JOIN patterns, and GROUP BY
idioms expected in analytical queries.  Given the same schema and question, it
consistently produces syntactically correct, executable SQL on the first attempt.

% ─────────────────────────────────────────────────────────────────────────────
\section{Query Comparison: Same Question, Two Models}
% ─────────────────────────────────────────────────────────────────────────────

To illustrate the difference concretely, we ran both models against the same questions
using the same database schema.  The relevant portion of the schema is shown below.

\begin{lstlisting}[style=sql, caption={Relevant schema (chicago\_crime table)}]
CREATE TABLE chicago_crime (
    "ID"                   BIGINT,
    "Case Number"          TEXT,
    "Date"                 TEXT,
    "Block"                TEXT,
    "IUCR"                 TEXT,
    "Primary Type"         TEXT,
    "Description"          TEXT,
    "Location Description" TEXT,
    "Arrest"               BOOLEAN,
    "Domestic"             BOOLEAN,
    "Beat"                 INTEGER,
    "District"             INTEGER,
    "Ward"                 DOUBLE PRECISION,
    "Community Area"       DOUBLE PRECISION,
    "FBI Code"             TEXT,
    "X Coordinate"         DOUBLE PRECISION,
    "Y Coordinate"         DOUBLE PRECISION,
    "Year"                 INTEGER,
    "Updated On"           TEXT,
    "Latitude"             DOUBLE PRECISION,
    "Longitude"            DOUBLE PRECISION,
    "Location"             TEXT,
    crime_type             TEXT
);
\end{lstlisting}

\subsection{Question 1: ``Which year had the most number of crimes?''}

\textbf{llama3.1:8b} was given this system prompt and bound to the
\texttt{sql\_db\_query} tool via the LangGraph \texttt{generate\_query} node:

\begin{lstlisting}[style=prompt, caption={System prompt delivered to llama3.1:8b}]
You are a PostgreSQL SQL expert.
Write a precise SELECT query for the question and call sql_db_query.
Use only column names that exist in the provided schema.
Rules:
- Enclose ALL column names in double quotes exactly as they appear in the schema.
- Use single quotes for string literals.
- Limit results to at most 3 rows unless the user asks otherwise.
- NEVER issue DML statements (INSERT, UPDATE, DELETE, DROP, etc.).
\end{lstlisting}

\noindent\textbf{Query generated by llama3.1:8b:}

\begin{lstlisting}[style=sql, caption={llama3.1:8b output for Question 1}]
SELECT year, COUNT(*) AS crime_count
FROM chicago_crime
GROUP BY year
ORDER BY crime_count DESC
LIMIT 1;
\end{lstlisting}

\noindent\textit{Issue:} The identifier \texttt{year} is left unquoted.  PostgreSQL
treats \texttt{year} as a reserved word in some dialects and, more critically, the
column is stored as \texttt{"Year"} (capital Y) in this case-sensitive database.
This query raises \texttt{column "year" does not exist}.

\vspace{0.8em}
\noindent\textbf{SQLCoder} received the full DDL inside the structured prompt template:

\begin{lstlisting}[style=prompt, caption={SQLCoder prompt template (abbreviated)}]
### Task
Generate a SQL query to answer [QUESTION]Which year had the most number
of crimes?[/QUESTION]

### Database Schema
<full CREATE TABLE DDL injected here>

### Answer
Given the database schema, here is the SQL query that answers
[QUESTION]Which year had the most number of crimes?[/QUESTION]
[SQL]
\end{lstlisting}

\noindent\textbf{Query generated by SQLCoder:}

\begin{lstlisting}[style=sql, caption={SQLCoder output for Question 1}]
SELECT "Year", COUNT(*) AS crime_count
FROM chicago_crime
GROUP BY "Year"
ORDER BY crime_count DESC
LIMIT 1;
\end{lstlisting}

\noindent\textit{Result:} Correct PostgreSQL quoting.  Executes successfully and
returns the year with the highest crime count.

\vspace{1em}

\subsection{Question 2: ``What are the top 10 crime types by count?''}

\noindent\textbf{Query generated by llama3.1:8b:}

\begin{lstlisting}[style=sql, caption={llama3.1:8b output for Question 2}]
SELECT "Primary Type", COUNT(*) AS crime_count
FROM chicago_crime
GROUP BY "Primary Type"
ORDER BY crime_count DESC
LIMIT 10;
\end{lstlisting}

\noindent\textit{Result:} Correct in this instance.  The model properly quoted
\texttt{"Primary Type"} here, but this was inconsistent across runs (the fallback
check-query node corrected earlier failures in the same session).

\vspace{0.8em}
\noindent\textbf{Query generated by SQLCoder:}

\begin{lstlisting}[style=sql, caption={SQLCoder output for Question 2}]
SELECT "Primary Type", COUNT(*) AS num_crimes
FROM chicago_crime
GROUP BY "Primary Type"
ORDER BY num_crimes DESC
LIMIT 10;
\end{lstlisting}

\noindent\textit{Result:} Consistently correct across multiple runs with no
fallback required.

\vspace{1em}

\subsection{Question 3: ``Which district has the highest number of arrests?''}

\noindent\textbf{Query generated by llama3.1:8b} (after check-query correction):

\begin{lstlisting}[style=sql, caption={llama3.1:8b output for Question 3 (post check-query)}]
SELECT "District", COUNT(*) AS arrest_count
FROM chicago_crime
WHERE "Arrest" = TRUE
GROUP BY "District"
ORDER BY arrest_count DESC
LIMIT 1;
\end{lstlisting}

\noindent\textbf{Query generated by SQLCoder:}

\begin{lstlisting}[style=sql, caption={SQLCoder output for Question 3}]
SELECT "District", COUNT(*) AS arrest_count
FROM chicago_crime
WHERE "Arrest" = true
GROUP BY "District"
ORDER BY arrest_count DESC
LIMIT 1;
\end{lstlisting}

\noindent Both models produced equivalent queries for this question.  The difference
here was in the number of pipeline steps required: \texttt{llama3.1:8b} needed the
check-query review node to correct an initial attempt that omitted the
\texttt{WHERE} clause, while SQLCoder generated the correct query on the first pass.

% ─────────────────────────────────────────────────────────────────────────────
\section{Key Findings}
% ─────────────────────────────────────────────────────────────────────────────

\begin{enumerate}[leftmargin=*, itemsep=4pt]
    \item \textbf{Fine-tuned beats general-purpose for SQL.}  SQLCoder's domain-specific
          training delivers consistently correct, executable SQL with no multi-step
          review cycle.  For a production text-to-SQL system, a fine-tuned model is
          the better choice.

    \item \textbf{RAG is unnecessary at this schema scale.}  Injecting the full DDL
          directly into the prompt was sufficient.  A RAG layer would have added
          infrastructure complexity with no accuracy benefit.

    \item \textbf{LangGraph adds robustness for weaker models.}  The multi-node
          graph (generate $\to$ check $\to$ execute) acts as a safety net for
          general-purpose models that produce imperfect SQL.  This overhead is
          unnecessary with SQLCoder.

    \item \textbf{Prompt engineering matters.}  The structured SQLCoder prompt template
          (\texttt{[QUESTION]…[/QUESTION]} / \texttt{[SQL]}) is critical for getting
          reliable completions.  Using a plain instruction-style prompt with SQLCoder
          degrades output quality.
\end{enumerate}

% ─────────────────────────────────────────────────────────────────────────────
\section{Next Steps}
% ─────────────────────────────────────────────────────────────────────────────

\subsection{Fine-tune SQLCoder on the Chicago Crime schema}

While SQLCoder performs well zero-shot, it occasionally produces queries that use
the wrong column alias or miss a domain-specific filter (e.g.\ filtering by
\texttt{crime\_type} vs.\ \texttt{"Primary Type"}).  Fine-tuning on a small set of
(question, SQL) pairs specific to the Chicago Crime database would further close this gap.

\subsection{Add a self-correction loop to the SQLCoder pipeline}

The LangGraph agent has a built-in check-query node that reviews SQL before execution.
A similar loop should be added to the SQLCoder pipeline: if \texttt{db.run(sql)}
raises an exception, feed the error message back to the model with a
\textit{``Fix this SQL error''} prompt and retry.

\subsection{Evaluate on a broader question set}

Current evaluation was qualitative and limited to a handful of questions.  A
structured benchmark with ground-truth SQL (e.g.\ Spider, BIRD, or a custom
Chicago Crime test set) would give a quantitative execution accuracy score and
expose failure modes more systematically.

\subsection{Expose as an API / chat interface}

Wrap the SQLCoder pipeline in a FastAPI endpoint or a Streamlit chat UI so that
non-technical analysts can query the dataset interactively without opening a
notebook.

\subsection{Support multi-table schemas}

The current database has a single denormalised table.  Extending the system to a
normalised schema (e.g.\ separating \texttt{district}, \texttt{location}, and
\texttt{crime} into linked tables) would test how well each model handles JOIN
queries and foreign-key relationships.

\subsection{Streaming and latency optimisation}

SQLCoder at 8B parameters on a local Ollama instance has a noticeable generation
latency, especially with a long DDL in the prompt.  Options to explore:
\begin{itemize}[leftmargin=*, itemsep=2pt]
    \item Quantised (Q4/Q5) model variants to reduce memory and latency.
    \item Schema caching: keep the tokenised DDL prefix cached in KV-cache (prefix
          caching, supported by vLLM).
    \item Moving to a larger, hosted SQLCoder variant (e.g.\ \texttt{sqlcoder-70b})
          for higher-accuracy production use.
\end{itemize}

% ─────────────────────────────────────────────────────────────────────────────
\section{Conclusion}
% ─────────────────────────────────────────────────────────────────────────────

This project demonstrated that a specialised, fine-tuned SQL generation model
(SQLCoder) outperforms a general-purpose model (llama3.1:8b) in a text-to-SQL
setting, even at the same parameter count.  Crucially, the complexity we initially
anticipated --- a full RAG pipeline with vector retrieval --- turned out to be
unnecessary.  Injecting the compact schema directly into the prompt was both simpler
and more reliable.

The system successfully answers natural-language questions about the Chicago Crime
dataset with no manual SQL writing, demonstrating that production-quality
text-to-SQL is achievable with open-source local models on modest hardware.

\end{document}
